\section{Experiments}
\label{sec:experiments}

\subsection{Experimental setup}

We evaluate \method{} under five aggregate-supervision themes:
\begin{enumerate}
  \item integer-valued sum labels;
  \item signed count labels with cancellation;
  \item Bernoulli count labels as a recovery of prior count loss;
  \item classwise count labels connecting to LLP;
  \item runtime of exact aggregate backends.
\end{enumerate}
Unless otherwise stated, MNIST-family experiments use LeNet-style instance networks, bags are sampled on the fly, and tables report mean $\pm$ sample standard deviation over random seeds when multiple seeds are available.
Hidden instance labels are never used for training; they are used only for evaluation.
Implementation details are summarized in Appendix~\ref{app:implementation-details}.
Each subsection changes the observed aggregate label while keeping the same finite-support convolutional likelihood layer.

\paragraph{Metrics.}
Across the experimental tables, aggregate MAE columns report the mean absolute difference between the predicted aggregate and the observed aggregate on test bags.
For each seed and configuration, all methods are evaluated on the same deterministic held-out bag sampler; train and test bags are drawn from disjoint underlying dataset splits.
For expected-value baselines, the predicted aggregate is the scalar expectation trained by the baseline.
For FS-Conv rows, the main aggregate prediction is the mode of the learned aggregate PMF; we also report expected-aggregate MAE for the main FS-Conv settings where this statistic is directly comparable to expected-value baselines.
For classwise count objectives, count MAE is the mean absolute class-count error averaged over classes and test bags,
\begin{equation}
  \frac{1}{|\mathcal{B}_{\rm test}|}
  \sum_{B\in\mathcal{B}_{\rm test}}
  \frac{1}{C}\sum_{c=1}^C |\hat{h}_{Bc}-h_{Bc}|.
\end{equation}
Rows labeled MSE denote MSE training objectives; MAE and accuracy columns are evaluation metrics computed after training.
Columns labeled instance accuracy use held-out instance labels only for evaluation, never for training.
Instance accuracy is never used for model selection unless explicitly stated.

\subsection{Integer-valued sum labels}

Digit-sum supervision is an integer-valued aggregate label: the learner observes a single scalar while each hidden instance label has ten possible values.
Here $z_i$ is the original MNIST digit value in $\{0,\ldots,9\}$ and the aggregate label is $Y_{\rm sum}=\sum_i z_i$.
Exact finite-support likelihood uses the full induced distribution over possible sums, whereas an expected-value baseline regresses only $\sum_i \mathbb{E}[z_i]$ to the observed sum.

\begin{table}[t]
\centering
\small
\caption{Integer-valued sum labels.  FS-Conv uses the full convolved PMF over sums, whereas MSE regresses only the expected sum.  MNIST FS-Conv rows report mean $\pm$ standard deviation over three seeds; SVHN rows use ResNet-18 backbones and are single-seed diagnostics.}
\label{tab:digit-sum}
\begin{tabular}{llcccc}
\toprule
Data & $n$/train bags & FS-Conv MAE & FS-Conv inst. acc. & MSE MAE & MSE inst. acc. \\
\midrule
MNIST & $10/1000$ & $0.612{\pm}.008$ & $.982{\pm}.000$ & $2.535$ & $.226$ \\
MNIST & $10/5000$ & $0.277{\pm}.026$ & $.992{\pm}.001$ & $1.023$ & $.746$ \\
MNIST & $50/1000$ & $4.901{\pm}.973$ & $.693{\pm}.228$ & $6.274$ & $.248$ \\
MNIST & $50/5000$ & $2.229{\pm}.616$ & $.887{\pm}.102$ & $2.594$ & $.817$ \\
SVHN scratch & $10/5000$ & $5.812$ & $.229$ & $5.328$ & $.249$ \\
SVHN pretr. & $10/5000$ & $1.324$ & $.955$ & $2.109$ & $.826$ \\
\bottomrule
\end{tabular}
\end{table}

For $n=10$ and $1000$ training bags, the MSE baseline has much worse aggregate error and hidden digit recovery than FS-Conv.
With more training bags it improves, but FS-Conv remains better in both aggregate MAE and hidden digit accuracy.
For $n=50$, the task is harder for both objectives because many digit assignments can give similar sums; FS-Conv still has lower aggregate error, while the lower hidden accuracy shows that large scalar-sum bags are underdetermined.
SVHN supports the same conclusion beyond MNIST when the visual backbone is strong: pretraining makes natural RGB house-number digit sums usable, and FS-Conv outperforms expected-sum regression with the same pretrained ResNet-18 backbone.
On SVHN with $5000$ bags, the pretrained FS-Conv model reaches $95.46\%$ hidden digit accuracy and $65.00\%$ exact sum accuracy, while the pretrained MSE baseline reaches $82.58\%$ hidden digit accuracy and $22.58\%$ exact sum accuracy.
For the MNIST FS-Conv rows, expected-aggregate MAE is $0.704{\pm}.029$, $0.306{\pm}.022$, $4.864{\pm}1.003$, and $2.183{\pm}.544$ for the four rows in Table~\ref{tab:digit-sum}, respectively.

\subsection{Signed aggregates with cancellation}

Signed counts test an aggregate that is not reducible to ordinary class proportions.
We again use the digit-$9$ target indicator, but each instance carries a known random sign; target-digit images contribute $+1$ or $-1$, while all other digits contribute $0$.
Even with random signs, the observed signed sum is many-to-one: positive and negative contributions can cancel, so different latent assignments can yield the same aggregate label.
FS-Conv models this full signed aggregate distribution, whereas the MSE baseline only matches the expected signed sum.

\begin{table}[t]
\centering
\scriptsize
\setlength{\tabcolsep}{3.5pt}
\caption{Signed aggregates with random known signs.  MSE regresses the expected signed sum $\sum_i s_i p_i$, while FS-Conv uses the full finite-support signed aggregate distribution.  MNIST FS-Conv rows report mean $\pm$ standard deviation over three seeds; FashionMNIST rows report seed 0--1 means; SVHN and CIFAR-10 rows are seed-0 pretrained ResNet-18 experiments.  A dash indicates that the corresponding paired MSE baseline was not evaluated.}
\label{tab:signed}
\begin{tabular}{llcccc}
\toprule
Data & $n$/train bags & MSE MAE & MSE inst. acc. & FS-Conv MAE & FS-Conv inst. acc. \\
\midrule
MNIST & $10/1000$ & $0.269$ & $.959$ & $0.095{\pm}.010$ & $.990{\pm}.001$ \\
MNIST & $10/5000$ & $0.047$ & $.995$ & $0.045{\pm}.010$ & $.996{\pm}.001$ \\
MNIST & $50/1000$ & $0.398$ & $.991$ & $0.243{\pm}.053$ & $.994{\pm}.001$ \\
MNIST & $50/5000$ & $0.203$ & $.995$ & $0.198{\pm}.022$ & $.996{\pm}.000$ \\
\midrule
FashionMNIST & $n=10$ & $0.121$ & $.987$ & $0.113$ & $.988$ \\
FashionMNIST & $n=50$ & $0.413$ & $.988$ & $0.381$ & $.990$ \\
\midrule
SVHN pretr. & $n=16$ & $0.745$ & $.926$ & $0.684$ & $.934$ \\
SVHN pretr. & $n=64$ & $1.639$ & $.914$ & $1.407$ & $.943$ \\
CIFAR-10 pretr. & $n=16$ & -- & -- & $0.343$ & $.974$ \\
CIFAR-10 pretr. & $n=64$ & -- & -- & $0.753$ & $.980$ \\
\bottomrule
\end{tabular}
\end{table}

Both methods recover accurate target detectors with enough data, but the likelihood uses the full signed aggregate distribution rather than only $\sum_i s_i p_i$ and gives lower signed-count error.
FashionMNIST shows the same pattern beyond MNIST: MSE is close in instance recovery, but FS-Conv is consistently better on the signed aggregate.
On SVHN with ImageNet-pretrained ResNet-18, the advantage is modest for $n=16$ but clearer for $n=64$, where FS-Conv improves instance accuracy from $91.37\%$ to $94.33\%$ and signed MAE from $1.639$ to $1.407$.
The CIFAR-10 pretrained rows evaluate whether the signed finite-support objective remains usable in a harder natural-image setting; they are not paired with MSE baselines.
They reach instance AUCs of $98.23\%$ for $n=16$ and $99.19\%$ for $n=64$.
This setting illustrates a distinction from nonnegative count supervision: cancellation changes the inverse problem even when the latent instance variable remains Bernoulli.
For the MNIST FS-Conv rows, posterior-mean signed MAE is $0.109{\pm}.011$, $0.050{\pm}.006$, $0.273{\pm}.054$, and $0.217{\pm}.014$ for the four rows in Table~\ref{tab:signed}, respectively.

\subsection{Bernoulli count likelihood as a special case}

Bernoulli count supervision validates the connection to prior count-loss work.
For MNIST-Bags, we make MNIST binary by choosing digit $9$ as the positive class and treating digits $0$--$8$ as negative.
The learner observes the number of target-digit images in the bag; hidden instance labels are used only for evaluation.
This case verifies that \method{} recovers the established Bernoulli count probability while sharing the same implementation pattern as the new aggregate types.
Standard attention MIL is not a like-for-like baseline here, because it is typically trained from a binary presence label $\mathbf{1}\{Y>0\}$, whereas this setting observes the full count $Y$.

\begin{table}[t]
\centering
\caption{Bernoulli count labels.  The Bernoulli count likelihood recovers accurate hidden target detectors on MNIST target-count bags.}
\label{tab:binary-countmil}
\begin{tabular}{lllcc}
\toprule
Data & $n$/train bags & Method & Count MAE & Instance acc. \\
\midrule
MNIST & $10/1000$ & exact count & $0.115$ & $0.988$ \\
MNIST & $10/5000$ & exact count & $0.090$ & $0.991$ \\
MNIST & $50/1000$ & exact count & $0.686$ & $0.987$ \\
MNIST & $50/5000$ & exact count & $0.460$ & $0.991$ \\
\bottomrule
\end{tabular}
\end{table}

This experiment is included as a recovery check rather than a novelty claim.
It confirms that the same likelihood layer used for signed counts and integer-valued sums also gives the standard Bernoulli count likelihood and accurate hidden target recovery.

\subsection{Classwise count objectives and LLP connections}

We include histogram supervision to position \method{} relative to recent learning-from-label-proportions objectives.
For MNIST histograms, the hidden category is the original digit class and the observed label is the vector of digit counts in the bag.
We evaluate a classwise composite likelihood: for each class, the model scores the observed one-vs-rest count with a Bernoulli count PMF.
This corresponds to the count-likelihood component of LLP-PVC, rather than a full joint multiclass histogram model.
The comparison here is intentionally limited: MSE matches expected class counts, while the classwise \method{} likelihood scores each observed count with the exact one-vs-rest count PMF.

\begin{table}[t]
\centering
\scriptsize
\setlength{\tabcolsep}{4pt}
\caption{Main classwise count-supervision results.  Count MAE is mean absolute class-count error.  \method{} denotes the same finite-support convolutional likelihood layer applied as a classwise composite likelihood over one-vs-rest counts, as described in Section~\ref{sec:convolutional-bag-prediction}.  A dash indicates that the corresponding paired MSE baseline was not evaluated.}
\label{tab:histogram-llp-pvc}
\begin{tabular}{lllcccc}
\toprule
Data & $n$/train bags & Backbone & MSE MAE & MSE inst. acc. & \method{} MAE & \method{} inst. acc. \\
\midrule
MNIST & $50/5000$ & LeNet & $1.684{\pm}.000$ & $.116{\pm}.000$ & $.080{\pm}.003$ & $.992{\pm}.000$ \\
CIFAR-10 & $16/1000$ & pretr. R18 & $.702$ & $.668$ & $.260$ & $.902$ \\
CIFAR-10 & $64/1000$ & pretr. R18 & $1.855$ & $.272$ & $.678$ & $.914$ \\
CIFAR-10 & $64/250$ & pretr. R18 & -- & -- & $1.115$ & $.803$ \\
\bottomrule
\end{tabular}
\end{table}

% Full histogram table retained for future use:
% MNIST & $10/1000$ & LeNet & MSE & $.181$ & $.955$ \\
% MNIST & $10/1000$ & LeNet & one-vs-rest count & $.032$ & $.983$ \\
% MNIST & $10/5000$ & LeNet & MSE & $.129$ & $.974$ \\
% MNIST & $10/5000$ & LeNet & one-vs-rest count & $.016$ & $.992$ \\
% MNIST & $50/1000$ & LeNet & MSE & $1.683$ & $.116$ \\
% MNIST & $50/1000$ & LeNet & one-vs-rest count & $.137$ & $.985$ \\
% MNIST & $50/5000$ & LeNet & MSE & $1.684$ & $.116$ \\
% MNIST & $50/5000$ & LeNet & one-vs-rest count & $.080$ & $.992$ \\
% CIFAR-10 & $16/1000$ & pretr. R18 & MSE & $.702$ & $.668$ \\
% CIFAR-10 & $16/1000$ & pretr. R18 & one-vs-rest count & $.260$ & $.902$ \\
% CIFAR-10 & $64/1000$ & pretr. R18 & MSE & $1.855$ & $.272$ \\
% CIFAR-10 & $64/1000$ & pretr. R18 & one-vs-rest count & $.678$ & $.914$ \\
% CIFAR-10 & $64/250$ & pretr. R18 & one-vs-rest count & $1.115$ & $.803$ \\

On the MNIST large-bag setting, classwise \method{} improves over MSE proportion matching in count MAE and hidden digit accuracy.
This supports the connection to class-count LLP objectives without making the histogram setting the main novelty claim.
The CIFAR-10 rows are narrower: they show that classwise \method{} is effective with ImageNet-pretrained ResNet-18 features and outperforms a pretrained MSE histogram baseline at $n=16$ and $n=64$.
However, the $n=64$, train $1000$ setting exposes four times as many image instances per epoch as $n=16$, train $1000$.
The fixed-instance-budget setting, $n=64$ with $250$ train bags, reaches only about $80.3\%$ final instance accuracy.
Thus the CIFAR-10 results should be read as evidence that \method{} can use aggregate-supervised image volume effectively under a strong representation, not as evidence that larger bags are intrinsically superior.

\subsection{Backend runtime}

Our FFT-tree backend computes batched multiclass one-vs-rest count PMFs on GPU and substantially reduces wall-clock time relative to a CPU dynamic-programming implementation in this benchmark.
This addresses a practical limitation of DP-based count losses: exact count-probability computation need not be tied to sequential CPU dynamic programming.
The benchmark in Table~\ref{tab:runtime} isolates aggregate PMF evaluation for batch size $256$ and $C=10$ classes, so each timing evaluates $2560$ binary count PMFs.
It measures backend speed rather than end-to-end training time or representation quality.

\begin{table}[t]
\centering
\caption{Runtime for multiclass one-vs-rest count PMFs with batch size $256$ and $C=10$ classes.  CPU DP is compared to GPU FFT-tree convolution for forward aggregate PMF computation only.}
\label{tab:runtime}
\begin{tabular}{rrrrrr}
\toprule
Bag size & CPU DP (s) & GPU FFT-tree (s) & Speedup & GPU mem. & Max error \\
\midrule
$64$ & $8.246$ & $0.00097$ & $8.5{\times}10^3$ & $13.8$ MB & $3.4{\times}10^{-7}$ \\
$128$ & $17.230$ & $0.00097$ & $1.8{\times}10^4$ & $28.6$ MB & $4.2{\times}10^{-7}$ \\
$256$ & $36.785$ & $0.00121$ & $3.0{\times}10^4$ & $55.6$ MB & $5.1{\times}10^{-7}$ \\
$512$ & $54.040$ & $0.00134$ & $4.0{\times}10^4$ & $113.3$ MB & $6.6{\times}10^{-7}$ \\
\bottomrule
\end{tabular}
\end{table}
